% ============================================================
%  Standalone Literature Review
%  Topic: GNN-Based Fault Detection and Localization in Power Grids
% ============================================================
\documentclass[12pt,a4paper]{article}

\usepackage[T1]{fontenc}
\usepackage[utf8]{inputenc}
\usepackage{amsmath,amssymb}
\usepackage{booktabs}
\usepackage{array}
\usepackage{geometry}
\usepackage{setspace}
\usepackage{cite}
\usepackage{hyperref}
\usepackage{microtype}

\geometry{margin=1in}
\onehalfspacing

\hypersetup{
    colorlinks=true,
    linkcolor=blue,
    citecolor=blue,
    urlcolor=blue,
    pdftitle={Literature Review on GNN-Based Fault Detection and Localization in Power Grids},
    pdfauthor={Anik Kumar Basu, Sabrina Shonda Snaha, Md. Sadik Mahmud Adive}
}

\title{\textbf{Literature Review on Graph Neural Networks for Fault Detection and Localization in Power Grids}}
\author{
Anik Kumar Basu \and
Sabrina Shonda Snaha \and
Md. Sadik Mahmud Adive\\
Department of Computer Science and Engineering\\
Bangladesh University of Business and Technology, Bangladesh
}
\date{June 2026}

\begin{document}

\maketitle

\begin{abstract}
Fault detection and localization are essential for secure and reliable smart-grid operation. As power systems integrate renewable energy sources, bidirectional power flows, distributed generation, and high-frequency sensing devices such as phasor measurement units (PMUs), conventional protection and monitoring methods face increasing difficulty. This literature review surveys the main research directions for power-system fault analysis, including traditional relay-based methods, shallow machine learning, deep learning, graph neural networks (GNNs), and spatiotemporal learning. The reviewed literature shows that topology-agnostic methods often fail to represent the physical structure of power networks, while static graph methods can underuse the temporal evolution of fault transients. Spatiotemporal GNNs are therefore a promising direction because they combine graph-based spatial learning with temporal sequence modeling. The review closes by identifying the main gaps around bus-level localization, robustness, real-data validation, and multi-task learning.
\end{abstract}

\noindent\textbf{Keywords:} literature review, graph neural networks, smart grid, fault detection, fault localization, PMU, GCN, LSTM, spatiotemporal learning.

% ============================================================
\section{Introduction}
% ============================================================

Modern power grids are becoming more complex because of renewable energy integration, changing load behavior, distributed energy resources, and high-speed measurement systems. These changes make fault detection and localization more challenging. A protection system must not only detect whether a fault has occurred, but also identify where the fault originated so that corrective action can be taken quickly.

The literature on this problem has evolved through several stages. Early approaches used impedance-based relays, threshold logic, and signal-processing techniques. Later, researchers introduced shallow machine-learning models such as support vector machines (SVMs), random forests, and artificial neural networks (ANNs). More recent studies use deep learning, including convolutional neural networks (CNNs), recurrent neural networks, and graph neural networks (GNNs).

The central challenge is that power-system faults are both spatial and temporal. They are spatial because disturbances propagate through buses, lines, transformers, and electrical coupling. They are temporal because voltage, current, frequency, and power signals evolve rapidly during pre-fault, fault, and post-fault periods. A strong fault-localization model should therefore represent both the grid topology and the transient signal behavior.

This review presents the main research directions and explains why spatiotemporal GNNs are a strong candidate for real-time fault detection and localization in power grids.

% ============================================================
\section{Traditional Fault Detection and Localization}
% ============================================================

Early fault-detection pipelines relied on impedance relays, threshold rules, and related signal-processing logic. These methods are fast and interpretable, but their accuracy depends on accurate line parameters, relay settings, and stable operating assumptions. In renewable-rich grids, these assumptions are less reliable because inverter-based resources, topology changes, and bidirectional power flow can alter fault signatures.

Liao et al. \cite{liao2021} discuss how smart-grid operation requires methods that can process large volumes of measurement data and adapt to changing system behavior. Chen et al. \cite{chen2020} similarly emphasize that data-driven fault diagnosis can address limitations of purely rule-based systems. Traditional methods remain important, but they are often insufficient for flexible, high-dimensional, PMU-based fault localization.

% ============================================================
\section{Shallow Machine-Learning Approaches}
% ============================================================

Shallow machine-learning methods were introduced to improve adaptability. Common examples include SVMs, decision trees, random forests, k-nearest neighbors, and shallow ANNs. These methods can learn nonlinear decision boundaries from measured voltage, current, power, or frequency features.

The main advantage of shallow learning is that it can outperform simple threshold rules when the training data are representative. For example, a classifier can learn patterns associated with different fault types or approximate fault locations. These models are also relatively easy to train compared with large neural networks.

The major limitation is feature dependence. Shallow models usually require manually designed features such as RMS values, harmonic components, sequence components, wavelet features, or statistical summaries. If the operating condition changes, those handcrafted features may no longer separate fault classes reliably. In addition, shallow models usually treat bus measurements as ordinary feature vectors and do not explicitly represent the physical grid topology.

For bus-level fault localization, ignoring topology is a serious limitation. Two buses may be electrically close even if their feature columns are far apart in a vector representation. This motivates graph-based learning methods that directly encode the power-network structure.

% ============================================================
\section{Deep Learning for Fault Diagnosis}
% ============================================================

CNN-based approaches reduce feature engineering by learning directly from waveform or feature-map inputs, but they still assume an ordered Euclidean layout and can be sensitive to bus indexing \cite{owerko2020}. Recurrent models, especially LSTMs, capture temporal evolution effectively \cite{hochreiter1997}, which makes them useful for fault transients. However, a pure temporal model does not explicitly represent electrical connectivity, so it can miss the spatial relationships that matter for bus-level localization.

% ============================================================
\section{Graph Neural Networks in Power Systems}
% ============================================================

GNNs address this mismatch by learning directly on the bus-line graph. Kipf and Welling \cite{kipf2017} provided the graph-convolutional foundation used in many later power-system studies. Chen et al. \cite{chen2020}, Donon et al. \cite{donon2019}, and James et al. \cite{james2020} showed that topology-aware learning improves diagnosis and localization performance in power networks. These results are important because they suggest that the grid topology itself should be part of the model, not just the input feature set.

% ============================================================
\section{Spatiotemporal Learning}
% ============================================================

Spatiotemporal graph models combine graph-based spatial aggregation with temporal sequence modeling. Yu et al. \cite{yu2018} demonstrated this idea in sequence prediction by learning both network structure and temporal dynamics. For fault localization, the same principle is especially relevant because the most discriminative evidence often appears in the first few cycles after fault inception. Still, many existing studies remain snapshot-based or focus on classification rather than bus-level localization. Robustness under measurement noise, topology changes, and class imbalance is also reported inconsistently. These gaps motivate the ST-GNN design used in this paper, which combines graph message passing with temporal modeling for joint fault-location and fault-type prediction.

% ============================================================
\section{Comparison of Method Families}
% ============================================================

Table \ref{tab:comparison} summarizes the major method families reviewed in this document.

\begin{table}[h]
\centering
\small
\caption{Comparison of Method Families for Power-Grid Fault Analysis}
\label{tab:comparison}
\begin{tabular}{@{}p{0.22\textwidth}p{0.30\textwidth}p{0.34\textwidth}@{}}
\toprule
\textbf{Method Family} & \textbf{Strength} & \textbf{Limitation} \\
\midrule
Relay and threshold methods & Fast, interpretable, practical for protection systems & Sensitive to parameter uncertainty and changing grid conditions \\
Shallow machine learning & Learns nonlinear patterns from measured data & Requires handcrafted features and often ignores topology \\
CNN-based deep learning & Learns features automatically from waveform or matrix inputs & Assumes Euclidean structure and depends on bus ordering \\
LSTM/recurrent models & Captures temporal evolution of fault signals & Does not explicitly model physical grid connectivity \\
GNN/GCN models & Uses grid topology directly through message passing & Often applied to static snapshots without temporal context \\
Spatiotemporal GNNs & Models both grid topology and time-varying transients & Requires careful dataset design and robustness validation \\
\bottomrule
\end{tabular}
\end{table}

% ============================================================
\section{Research Gap}
% ============================================================

The reviewed literature reveals several research gaps. First, many studies emphasize fault detection or fault-type classification, while bus-level localization remains more difficult. Second, topology-agnostic models may perform well on fixed datasets but lack physical consistency under bus reordering or topology changes. Third, spatial GNN models often underuse temporal fault signatures. Fourth, many studies report accuracy but provide limited robustness analysis under measurement noise, class imbalance, and N-1 contingencies.

These gaps suggest the need for a model that jointly learns topology and time. A spatiotemporal GNN can address this by using graph convolution for spatial dependencies and sequence modeling for transient behavior. Multi-task learning can also help because fault type and fault location are related outputs.

% ============================================================
\section{Direction for Current Research}
% ============================================================

The reviewed works support a research direction based on three ideas. First, fault signals are graph-structured because physical connections determine how disturbances propagate. Second, fault signals are time-dependent because useful localization information appears during short transient windows. Third, fault location and fault type are related, so a shared representation can support both tasks.

For this reason, a model that combines GCN layers, LSTM layers, and multi-task output heads is well aligned with the limitations found in existing work. Synthetic data generation tools such as Pandapower \cite{thurner2018} can support systematic experiments by varying bus location, fault type, fault resistance, and loading level.

% ============================================================
\section{Conclusion}
% ============================================================

This literature review shows a clear progression from traditional protection methods to topology-aware spatiotemporal learning. Traditional and shallow-learning methods remain useful, but they struggle with high-dimensional PMU data, changing grid conditions, and bus-level localization. CNN and LSTM methods improve feature learning but do not fully solve the topology-time coupling problem. GNNs provide a natural way to represent power grids, while spatiotemporal GNNs extend this advantage by modeling transient dynamics.

Overall, the literature supports the use of spatiotemporal graph neural networks for real-time fault detection and localization. Future research should focus on real-data validation, robustness under noise and topology changes, explainability, and transfer learning across different grid systems.

% ============================================================
\begin{thebibliography}{9}

\bibitem{liao2021}
W.~Liao, B.~Bak-Jensen, J.~R.~Pillai \emph{et al.},
``Review of Graph Neural Networks for Smart Grids: An Overview,''
\emph{IEEE Power \& Energy Society General Meeting (PESGM)}, 2021.

\bibitem{owerko2020}
D.~Owerko, F.~Gama, and A.~Ribeiro,
``Optimal Power Flow using Graph Neural Networks,''
\emph{IEEE Int.\ Conf.\ Acoustics, Speech and Signal Processing (ICASSP)}, 2020.

\bibitem{chen2020}
K.~Chen, J.~Hu, and Y.~Zhang,
``Power System Fault Diagnosis Based on Graph Convolutional Networks,''
\emph{IEEE Int.\ Conf.\ Energy Internet and Energy System Integration (EI2)}, 2020.

\bibitem{donon2019}
B.~Donon, B.~Donnot, I.~Guyon, and M.~Sebag,
``Graph Neural Solver for Power Systems,''
\emph{Int.\ Joint Conf.\ Neural Networks (IJCNN)}, 2019.

\bibitem{yu2018}
B.~Yu, H.~Yin, and Z.~Zhu,
``Spatiotemporal Graph Convolutional Networks: A Deep Learning Framework
for Traffic Forecasting,''
\emph{IJCAI}, 2018.

\bibitem{james2020}
J.~James, Y.~C.~Chen, and A.~Dominguez-Garcia,
``Fault Location in Power Distribution Networks via Graph Convolutional Networks,''
\emph{IEEE North American Power Symposium (NAPS)}, 2020.

\bibitem{kipf2017}
T.~N.~Kipf and M.~Welling,
``Semi-Supervised Classification with Graph Convolutional Networks,''
\emph{Int.\ Conf.\ Learning Representations (ICLR)}, 2017.

\bibitem{hochreiter1997}
S.~Hochreiter and J.~Schmidhuber,
``Long Short-Term Memory,''
\emph{Neural Computation}, vol.~9, no.~8, pp.~1735--1780, 1997.

\bibitem{thurner2018}
L.~Thurner \emph{et al.},
``Pandapower---An Open-Source Python Tool for Convenient Modeling, Analysis,
and Optimization of Electric Power Systems,''
\emph{IEEE Trans.\ Power Systems}, vol.~33, no.~6, pp.~6510--6521, 2018.

\end{thebibliography}

\end{document}
